\section{Proof Sketches and Technical Notes}
\label{app:proofs}

\subsection{Privacy-minimality
(Section~\ref{sec:sacrifice_model})}

\paragraph{Theorem~\ref{thm:privacy} + Proposition~\ref{prop:privacy_institutional}
strategy.}
The three formal properties (P1, P3, P4) are proved
independently from the channel's construction; the two
institutional preconditions (P2, P5) are documented separately
in Proposition~\ref{prop:privacy_institutional}.  Proof
strategy:
P1~(no interior access): by construction of the observation
tuple.
P3~(discretization): by bounding channel bandwidth with trade
frequency.
P4~(commitment composability): by composition
with~\citep[Definition~6]{lasser2026exo}.
P2 and P5 are not theorem content: P2's zero-contribution
property is formal (a consequence of the aggregate's summation
structure), but the consent interpretation depends on
institutional opt-in conventions; P5 depends on public-ledger
infrastructure external to the paper's formal scope, and is
additionally split between money-channel events (observable
from public ledger) and time-channel events (require an
additional public wage schedule or imputation rule).

\subsection{Aggregate lower bound
(Section~\ref{sec:lower_bound})}

\paragraph{Theorem~\ref{thm:aggregate_bound} strategy.}

\emph{Part~A (bundle-level, S0--S2 and S4(a),
poset-independent bundles, genuine-exchange conditions):}
\begin{enumerate}
\item Revealed-preference observation: $\Ui[_n]((H_n \setminus
  X_n) \cup Y_n) \geq \Ui[_n](H_n)$.
\item Cancellation to bundle comparison (S4(a) +
  Definition~\ref{def:sacrifice_event} conditions (i)--(iv),
  via Lemma~\ref{lem:cancellation}):
  $\Ui[_n](Y_n) \geq \Ui[_n](X_n)$.
\item Calibration chain (S0 + S2), S0 stated directly in
  window-integrated form:
  $\volR^{[W]}(Y_n) \geq \Ui[_{n}](Y_n) \geq \Ui[_{n}](X_n)
  \geq \Delta\volL(X_n)$.  S0's consequent is against
  $\volR^{[W]}$ rather than against the instantaneous
  $\volR_{t_n}$: a sacrifice event records acquisition, not an
  instantaneous exercise witness, so the bound is about
  cumulative exercise of $Y_n$ anywhere in $W$
  (Definition~\ref{def:volR},
  Section~\ref{sec:assumptions}).
\item Poset-independence of bundles
  (Definition~\ref{def:category_partition}) gives M4 additivity
  (inheritance proved in Theorem~\ref{thm:axiom_inheritance}
  via sub-poset restriction):
  $\volR^{[W]}(\bigcup Y_n) = \sum_n \volR^{[W]}(Y_n) \geq
  \sum_n \Delta\volL(X_n)$.
\end{enumerate}

\emph{Part~B (per-capability refinement, additionally S3--S5
and within-bundle poset-independence; conditional on S5, see
Remark~\ref{rem:s5_heavy}):}
\begin{enumerate}
\setcounter{enumi}{3}
\item Part~B is stated for events with $\Delta\volL(X_n) > 0$
  and bundles $Y_n$ satisfying within-bundle poset-independence
  (Definition~\ref{def:bundle_decomp},
  Remark~\ref{rem:within_bundle_indep}); bundles with internal
  cooperative activation fall back to Part~A only.  Positivity
  $\volR(Y_n) > 0$ then follows from Part~A.
\item Within-bundle poset-independence plus~M4 gives
  $\volR(Y_n) = \sum_{c \in Y_n} \volR(c)$, making the
  partition-of-unity constraint
  $\sum_{c \in Y_n} \alpha_{n,c} = 1$ compatible with S5's
  share inequality $\alpha_{n,c} \leq \volR(c)/\volR(Y_n)$
  (equality achievable).  Combined: $\volR(c) \geq \alpha_{n,c}
  \cdot \Delta\volL(X_n)$.  S3 (bundle coherence) ensures the
  hedonic decomposition is non-degenerate so $\alpha_{n,c}$ is
  well-specified.
\item Per-capability max-attribution across overlapping events
  (Proposition~\ref{prop:monotone}).
\item Summation over capabilities in $U$ (pairwise
  poset-independent) via a second M4 application.
\end{enumerate}

Key technical challenges:
\begin{itemize}
\item \emph{Poset-independence (Part~A):} the category
  partition $\mathcal{C}$
  (Definition~\ref{def:category_partition}) must be a genuine
  poset partition.  Set-disjointness is necessary but not
  sufficient (Remark~\ref{rem:poset_indep}).
\item \emph{Decomposition validity (S5, Part~B):} hedonic
  regression~\citep{rosen1974hedonic} produces attribution
  weights, not $\volR$ share lower bounds.  S5 is the explicit
  bridge.
\item \emph{Calibration (S0):} the weakest assumption in the
  chain but the hardest to verify empirically (Open
  Question~\ref{oq:calibration}).
\end{itemize}

\subsection{Axiom inheritance
(Section~\ref{sec:axiom_inheritance})}

\paragraph{Theorem~\ref{thm:axiom_inheritance} strategy.}
Axiom-by-axiom analysis:
M1--M5 are direct from properties of summation/max over
non-negative terms (detailed in the body; M4 is proved
constructively via sub-poset restriction).
M6 for $\volR$ as a measure on a fixed exercised sub-poset is
inherited from $\volL$ by restriction---no axiom failure
there.  The paper's non-self-balancing claim rests instead on
\emph{capability-removal non-invariance}: removing an
unexercised capability leaves $\volR$ unchanged while $\volL$
drops by the removed weight
(Example~\ref{ex:volR_removal_invariance}).  This is a
system-level property (of the mapping
$\mathcal{P} \mapsto \volR^{\mathcal{P}}$), not a measure-axiom
failure.  M6(b) for $\volRlower$ is witnessed separately by
Example~\ref{ex:m6_fail_volRlower}: max-attribution
subadditivity.  Corollary~\ref{cor:not_self_balancing} follows
from~(a)'s non-invariance, not from~(b)---a lower bound failing
measure axioms is unremarkable; what breaks self-balancing
transfer is that $\volR$ doesn't penalize shedding of unused
capabilities.
